% ============================================================================
%  Chapter 7. Evaluation: Four Independent Measurements of Failure
%
%  \input file. Requires booktabs, amsmath, natbib in the parent preamble.
%  Bibliography: accord/thesis/refs.bib
%
%  Traceability: the audit scripts over the 299 pharmacist-annotated pairs,
%  transfer_results_correct.json, the MedNLI fine-tuning logs, and
%  accord/CORRECTIONS.md (C2, C4, C5).
% ============================================================================

\chapter{Evaluation: Four Independent Measurements of Failure}
\label{ch:eval}

\section{Evaluation Design}
\label{sec:eval-design}

A negative result is only worth reporting if it is difficult to escape. A single
unfavourable metric invites the response that the metric was the wrong one, the
threshold was badly chosen, the class balance was hostile, or the sample was too
small. This chapter is constructed to make those responses unavailable, by
measuring the same system four ways that fail for different reasons if they fail
at all.

\begin{enumerate}
\item \textbf{Ranking.} Does the model's contradiction probability order the
      pharmacist's labels better than chance? This is threshold-free and
      insensitive to class balance. Section~\ref{sec:m1}.
\item \textbf{Operating points.} Is there \emph{any} threshold at which the
      model is usable? This removes the objection that $0.5$ was the wrong
      cut. Section~\ref{sec:m2}.
\item \textbf{Severity.} Does the model's harm score correlate with the
      pharmacist's $1$ to $5$ danger rating? This asks a different question of a
      different output head. Section~\ref{sec:m3}.
\item \textbf{Comparison.} How does the model rank against a general-purpose
      language model, zero-shot and few-shot, and against the weak labeller it
      was distilled from, on the same $299$ pairs? This removes the objection
      that the task is simply impossible. Section~\ref{sec:m4}.
\end{enumerate}

All four are computed on the $299$ pharmacist-annotated pairs.

\paragraph{A correction that governs every number in this chapter.} Earlier
drafts of Chapters~\ref{ch:eval} to~\ref{ch:discussion} stated that the
configuration trained on the full $1{,}978$-pair pool, which contains these $299$
pairs, was not reported anywhere in this thesis. \textbf{That statement was
untrue.} Every point metric, ranking statistic, threshold sweep and severity
correlation in those drafts was computed from
\texttt{gold\_standard\_accord\_results.csv}, and that file is the
\emph{contaminated} run: it records $11$ true positives of $12$, a recall no
decontaminated seed reaches (\texttt{CORRECTIONS.md} C16). The figures reported
below replace it with the decontaminated configuration, retrained on the held-out
complement and evaluated over three seeds. The decontaminated run is worse than
the contaminated one on recall, precision, $F_1$ and $\kappa$, so the negative
conclusion this chapter draws is strengthened by the correction rather than
weakened by it. Quantities that have not been recomputed under decontamination
are marked $\dagger$ where they appear and are listed in
Section~\ref{sec:eval-limitations}.

Section~\ref{sec:eval-scoping} then addresses the obvious objection: given the
exemplar defect of Chapter~\ref{ch:candidates}, can any of this be believed? The
answer is a qualified yes, and the qualification is stated precisely rather than
waved at.

\section{A Component Sanity Check: BioMedBERT on MedNLI}
\label{sec:mednli}

\subsection{Purpose}

Before reporting that a system fails, it is worth establishing that the pieces
were assembled competently. MedNLI \citep{romanov2018mednli} is the standard
clinical natural language inference benchmark: premise and hypothesis pairs derived
from MIMIC-III, labelled entailment, contradiction, or neutral. It is the
closest well-established task to the one this thesis attempts, and it has an
established performance range.

\subsection{Result}

Table~\ref{tab:mednli} reports the two configurations.\footnote{Provenance
note. The four accuracy and macro-$F_1$ figures in Table~\ref{tab:mednli}
($83.90$, $83.91$, $72.93$, $72.74$) are carried forward from a run log that is
not among this project's surviving artifacts: no run log, no metrics file, and
no checkpoint for either configuration could be located. They are the only
numbers in this thesis whose producing artifact could not be found, and they are
retained rather than deleted so that a reader encountering them elsewhere knows
their status. The qualitative claim they support --- that the backbone,
tokenisation, training loop, and evaluation harness are sound, and that the LoRA
configuration carries a material accuracy cost --- does not carry any of this
chapter's conclusions, every one of which rests on artifacts that are present.}

\begin{table}[t]
\centering
\small
\begin{tabular}{lrrr}
\toprule
\textbf{Configuration} & \textbf{Trainable} & \textbf{Test acc.} & \textbf{Macro $F_1$} \\
\midrule
BioMedBERT, full fine-tune, 10 epochs        & full backbone & 83.90\% & 83.91\% \\
BioMedBERT $+$ LoRA ($r{=}16$, $\alpha{=}32$), 5 epochs & 592,131 & 72.93\% & 72.74\% \\
\midrule
Published range for comparable systems       & n/a & \multicolumn{2}{c}{83 to 87\%} \\
\bottomrule
\end{tabular}
\caption{MedNLI results. The full fine-tune sits inside the published range for
comparable encoder-based systems. This establishes that the backbone, the
tokenisation, the training loop, and the evaluation harness are sound. The LoRA
configuration used for ACCORD loses roughly eleven accuracy points, which is a
material cost that was not accounted for in the project's design.}
\label{tab:mednli}
\end{table}

\subsection{What this row is, and what it is not}

\textbf{Table~\ref{tab:mednli} is a BioMedBERT baseline. It is not ACCORD.} Zero
ACCORD components are present in either model: no cross-note encoder, no
authority conditioning, no prior conditioning, no harm estimation. It is a
sentence-pair classifier on a public benchmark.

This distinction is emphasised because the alternative framing, reporting
$83.90\%$ under the project's own system name, would be the single most
misleading move available in this thesis. A reader encountering
``ACCORD: $83.9\%$ on MedNLI'' followed by the results in
Section~\ref{sec:eval-point} would reasonably conclude that a strong model
failed on a hard task. What actually happened is that a strong \emph{backbone}
performs as expected on a benchmark it was designed for, and a differently
constructed system built on that backbone fails on a different task for reasons
that have nothing to do with the backbone's capability.

\subsection{One substantive finding}

The eleven-point gap between full fine-tuning and LoRA is worth recording. LoRA
is often described as approximately loss-free relative to full fine-tuning
\citep{hu2022lora}, and on many tasks it is. Here, at $r = 16$ with $592{,}131$
trainable parameters\footnote{Unlike the accuracies, this count is retained
because it derives exactly. LoRA at $r = 16$ on the query and value projections
of all twelve BERT-base layers contributes
$12 \times 2 \times (768 \times 16 + 16 \times 768) = 589{,}824$ parameters, and
a three-class classification head contributes $768 \times 3 + 3 = 2{,}307$, for a
total of $592{,}131$. It is independently corroborated by
\texttt{results/accord/stage0/config1\_baseline\_bottleneck\_head.json}, which
reports $110{,}074{,}887$ trainable parameters, that is, BERT-base's
$109{,}482{,}240$ plus a $592{,}647$-parameter head.} and a shorter schedule, it
is not. ACCORD used the LoRA
configuration. Whether the gap reflects rank, schedule length, or the specific
demands of clinical inference is untested here. [TODO: run a rank sweep
$r \in \{8, 16, 64, 256\}$ at matched epochs on MedNLI and report the
accuracy and parameter frontier. The two rows in Table~\ref{tab:mednli} differ in
both rank and schedule and therefore do not isolate either.]

\section{Point Metrics on the Gold Set}
\label{sec:eval-point}

\begin{table}[t]
\centering
\small
\begin{tabular}{lr}
\toprule
\textbf{Quantity} & \textbf{Value} \\
\midrule
Annotated pairs                                   & 299 \\
\quad true contradictions                         & 12 \\
\quad not contradictions                          & 241 \\
\quad ambiguous                                   & 46 \\
\midrule
Pairs ACCORD predicts \emph{contradiction} at $\tau = 0.5$ & 278$^{\dagger}$ \\
\midrule
Precision                                         & $0.0424 \pm 0.0023$ \\
Recall                                            & $0.7778 \pm 0.0393$ (TP $9$, $10$, $9$ of $12$) \\
Accuracy                                          & 0.115$^{\dagger}$ \\
Pooled base rate                                  & 0.040 (12/299) \\
\bottomrule
\end{tabular}
\caption{Point metrics at the default threshold. Precision and recall are the
decontaminated configuration's three-seed mean $\pm$ standard deviation; the
recall of $0.7778$ replaces the $0.917$ ($11/12$) of the contaminated run, whose
three decontaminated seeds recover $9$, $10$ and $9$ of the $12$ confirmed
contradictions. Rows marked $\dagger$ are contaminated-run quantities that have
not been recomputed under decontamination and are retained only to support the
denominator reconstruction below. The model remains a near-constant positive
predictor evaluated against a heavily negative population, and the precision of
$0.0424$ should be read against the base rate rather than in isolation.}
\label{tab:point}
\end{table}

\paragraph{On denominators.} Precision and accuracy in Table~\ref{tab:point} are
computed on the $253$ non-ambiguous pairs, while the predicted-positive count of
$278$ is over all $299$.\footnote{The reconstruction below is of the
\emph{contaminated} run's released file, which is the only configuration for
which the released per-pair decisions exist; it is retained because it is what
fixes the denominators, not because its cell counts still stand. The
threshold sweep in Section~\ref{sec:m2} reports precision $0.0474$ at a
positive rate of $1.0$, and $12/253 = 0.0474$, which fixes the denominator at
$253$. With $11$ true positives, a precision of $0.047$ requires $234$ predicted
positives among those $253$, and an accuracy of $0.115$ then follows as
$29/253$. The $278$ figure over all $299$ implies that $44$ of the $46$
ambiguous pairs are also predicted positive. [TODO: confirm these denominators
directly in the audit script and report them explicitly in the results file. The
present mixture of $n{=}253$ and $n{=}299$ denominators in a single table is a
reporting defect, corrected here by annotation rather than by recomputation.]}
On the $253$-pair basis the comparable base rate is $12/253 = 4.74\%$, and
ACCORD's decontaminated precision of $4.24\% \pm 0.23$ is \emph{below} it.

\paragraph{The reading.} A detector whose precision sits half a percentage point
below the base rate of the pool it is scoring is not a detector. It is a
relabelling of the candidate pool. The recall of $0.778 \pm 0.039$ is real, but
it is the recall that a near-constant positive predictor achieves: predicting
\emph{contradiction} for everything yields recall $1.0$ at precision equal to
the base rate, and ACCORD is a small perturbation of that predictor.

This is the point at which most evaluations stop, and it is also the point at
which the result is most easily dismissed. The next four sections make it
undismissable.

\section{Measurement 1: The Score Carries No Ranking Information}
\label{sec:m1}

\subsection{Why ranking is the decisive test}

Precision and recall depend on a threshold, and thresholds can be argued about.
Area under the ROC curve does not. It is the probability that a randomly drawn
positive receives a higher score than a randomly drawn negative. If the model
has learned \emph{anything} about the task, even something too weak to support a
usable threshold, AUROC exceeds $0.5$. If AUROC is $0.5$, the score is
statistically independent of the label.

\begin{table}[t]
\centering
\small
\begin{tabular}{llrr}
\toprule
\textbf{Score} & \textbf{Label treatment} & \textbf{$n$} & \textbf{AUROC} \\
\midrule
\texttt{accord\_contra\_prob} & all pairs, ambiguous excluded from positives & 299 & 0.5145 \\
\texttt{accord\_contra\_prob} & ambiguous pairs dropped                       & 253 & 0.5145 \\
\texttt{accord\_contra\_prob} & ambiguous counted as positive                 & 299 & 0.490$^{\dagger}$ \\
\midrule
\texttt{accord\_harm\_score}  & all pairs                                     & 299 & 0.504$^{\dagger}$ \\
\bottomrule
\end{tabular}
\caption{Ranking performance under three treatments of the ambiguous class. The
first two rows are the decontaminated three-seed ensemble, $0.5145$, whose
per-seed mean is $0.5211 \pm 0.0201$; they replace the $0.460$ of the
contaminated run, which reported the same value under both treatments. The
decontaminated recomputation exists only on the $253$-pair core, so the rows
marked $\dagger$ remain contaminated-run values and have not been recomputed. The
harm score, which Chapter~\ref{ch:arch} established was never trained, sits at
$0.504$: the value expected of a randomly initialised head.}
\label{tab:auroc}
\end{table}

\subsection{The result}

Table~\ref{tab:auroc} reports AUROC of $0.5145$ for the contradiction
probability, the three-seed ensemble of the decontaminated configuration, with a
per-seed mean of $0.5211 \pm 0.0201$. This is fractionally above chance, and the
supporting statistics establish that it is indistinguishable from chance:

\begin{itemize}
\item \textbf{Average precision} $0.0579 \pm 0.0034$ across the three
      decontaminated seeds, against a base rate of $0.040$. Average precision is
      the appropriate summary under this degree of class imbalance
      \citep{saito2015prcurve}, and it sits at the no-skill line.
\item \textbf{Mann-Whitney $U$} comparing the score distributions of positives
      and negatives: $p = 0.867$. The two distributions are not separated.
\item \textbf{Bootstrap $95\%$ confidence interval} on AUROC, $B = 10{,}000$
      resamples \citep{efron1994bootstrap}: $[0.321, 0.703]$. The interval
      contains $0.5$ comfortably. It is also very wide, which is the honest
      consequence of $12$ positives.
\end{itemize}

\paragraph{Neither configuration is certified subject-disjoint.} The $0.5145$
above is not a clean number, and the reason should be stated at the point of
first use rather than in the limitations. The decontaminated run's own
train/eval boundary was enforced at \emph{row} level: rows carrying the
\texttt{pharmacist\_overlap} flag were removed, but subjects were not. A direct
check (\texttt{CORRECTIONS.md} C18) finds that $158$ of the $266$ subjects
spanned by the evaluated pairs --- $59.4\%$ --- also appear in the
decontaminated training pool, under different \texttt{(drug, rule)}
combinations. So the contaminated configuration leaked at row level and the
decontaminated configuration leaks at subject level, and \textbf{neither is
certified subject-disjoint from its evaluation set}. The only subject-disjoint
measurement anywhere in this project is the plain control baseline of
\texttt{CORRECTIONS.md} C19, which excludes the $291$ overlapping rows and
asserts disjointness in code; it lands at AUROC $0.497 \pm 0.035$. The three
measurements --- contaminated $0.4595$, decontaminated $0.5145$, subject-disjoint
control $0.497$ --- agree, which is the useful outcome: leakage at either level
did not manufacture signal, because there was none to manufacture.

\subsection{The top-$k$ check}

AUROC is a global summary and can obscure behaviour at the extremes, which is
where a triage system would actually operate. \textbf{The $15$ highest-scored
pairs in the evaluation set contain one true contradiction.} Under the null
hypothesis of random ordering, the expected count is $15 \times 12/299 = 0.6$,
so observing one is unremarkable in itself. That is precisely the point: the
model's most confident predictions are indistinguishable from a random draw. The
count of one is the decontaminated ensemble's; the contaminated run reported
zero, and across the three decontaminated seeds the counts are $1$, $0$ and $0$.

For a deployed triage system this is the operationally decisive number. The
entire value proposition of such a system is that a clinician reviews the top of
the ranked list. Here the top of the list is noise.

\subsection{The harm score}

The harm score's AUROC of $0.504$ at $p = 0.96$ is reported for completeness and
should be interpreted as Chapter~\ref{ch:arch}, Section~\ref{sec:c3} instructs:
this head received no loss term at any point in training. Its output range
across the $299$ pairs is $0.398$ to $0.610$, the range of an untrained sigmoid.
A null result here is a measurement of an untrained module, not evidence about
whether clinical harm is predictable.

\section{Measurement 2: No Operating Point Exists}
\label{sec:m2}

\subsection{The objection this forecloses}

The natural response to Section~\ref{sec:eval-point} is that $\tau = 0.5$ is the
wrong threshold for a $4\%$-prevalence task and that a lower threshold, or a
prior correction, would produce a usable system. This section shows that no
threshold produces a usable system, because there is nothing in the score
ordering for a threshold to exploit.

\subsection{The sweep}

\begin{table}[t]
\centering
\small
\begin{tabular}{lrrr}
\toprule
\textbf{Threshold $\tau$} & \textbf{Precision} & \textbf{Recall} & \textbf{Positive rate} \\
\midrule
0.05 & 0.0474 & 1.000 & 1.000 \\
0.10 & 0.0474 & 1.000 & 1.000 \\
0.15 & 0.0474 & 1.000 & 1.000 \\
\bottomrule
\end{tabular}
\caption{Threshold sweep. At every threshold in the low range the model predicts
positive for \emph{every} pair. The precision of $0.0474$ is exactly $12/253$,
the base rate: at a positive rate of $1.0$ precision is arithmetically forced to
the base rate. The three rows are identical because no score in the evaluation
set falls below $0.15$. $\dagger$~This sweep is computed from the contaminated
run's released scores and has not been recomputed under decontamination.}
\label{tab:sweep}
\end{table}

\begin{table}[t]
\centering
\small
\begin{tabular}{lr}
\toprule
\textbf{Target recall} & \textbf{Achieved precision} \\
\midrule
0.500 & 5.4\% \\
0.750 & 4.1\% \\
0.917 & 4.1\% \\
\bottomrule
\end{tabular}
\caption{Precision at fixed recall. Even at $50\%$ recall, precision reaches
only $5.4\%$: roughly eighteen false alarms for every true contradiction. The
precision-recall curve is essentially flat at the base rate, which is the
signature of a score that does not order the labels. $\dagger$~The whole table is
computed from the contaminated run's released scores and has not been recomputed
under decontamination; the $0.917$ target-recall row is a grid point chosen
because it was that run's achieved recall, and the decontaminated configuration
achieves $0.778 \pm 0.039$. [TODO: recompute the curve on the decontaminated
scores, replace the $0.917$ grid point, and reconcile the denominators, since the
$4.1\%$ figure and the precision reported at $\tau = 0.5$ in
Table~\ref{tab:point} are computed on different bases ($n = 299$ against
$n = 253$).]}
\label{tab:pr}
\end{table}

\subsection{The reading}

Tables~\ref{tab:sweep} and~\ref{tab:pr} together establish that the model's
scores occupy a narrow band above $0.15$ and that within that band the ordering
is uninformative. No threshold, no prior correction, and no post-hoc calibration
recovers a usable operating point, because calibration reshapes a score
distribution and cannot create an ordering that is not there
\citep{guo2017calibration}.

This is the answer to a question Chapter~\ref{ch:weak} deferred. The prevalence
mismatch in the weak labels could in principle have been corrected after the
fact by shifting the decision threshold. It could not be corrected here, because
the damage was not a shifted threshold on an informative score. It was the loss
of the score's informativeness altogether.

\section{Measurement 3: Severity Ranking}
\label{sec:m3}

The pharmacist assigned a $1$ to $5$ danger rating to each pair. Two Spearman rank
correlations were computed against it:

\begin{table}[h]
\centering
\small
\begin{tabular}{lrrr}
\toprule
\textbf{Model score} & \textbf{Spearman $\rho$} & \textbf{$p$} & \textbf{$n$} \\
\midrule
\texttt{accord\_harm\_score}  & $-0.0001$ & 0.999 & 299 \\
\texttt{accord\_contra\_prob} & $-0.029$  & 0.62  & 299 \\
\bottomrule
\end{tabular}
\caption{Severity ranking. Both correlations are null. The harm score's
$\rho = -0.0001$ at $p = 0.999$ is about as close to an exact null as a finite
sample produces. $\dagger$~Both rows are computed from the contaminated run's
released scores and have not been recomputed under decontamination.}
\label{tab:severity}
\end{table}

\paragraph{How to read this, and how not to.} The harm-score row is \emph{not}
evidence that clinical severity is unpredictable. Chapter~\ref{ch:arch},
Section~\ref{sec:c3} established that the harm head receives no loss term, its
guideline bank is twenty random vectors, and its outputs span the range of an
untrained sigmoid. A correlation of $-0.0001$ between a randomly initialised
head's outputs and a clinical rating is the expected result and confirms the
architectural finding.

\textbf{The severity-ranking research question was therefore not tested by this
work.} It remains open. What this thesis establishes about it is narrower and
still useful: a severity head can be specified, implemented, run end to end,
included in a results table, and reported as a contribution without ever
appearing in the loss, and the only measurement that catches this is the
empirical range of its outputs.

The contradiction-probability row is a genuine, if unsurprising, null: a score
that does not separate contradictions from non-contradictions
(Section~\ref{sec:m1}) has no reason to rank them by danger.

\section{Measurement 4: System Comparison}
\label{sec:m4}

\subsection{The objection this forecloses}

A failure that is universal across methods is a statement about the task, not
about the system. If every approach performs at the base rate, the correct
conclusion is that cross-note contradiction detection at $4\%$ prevalence from a
two-sentence view is not solvable, and the specific defects catalogued in
Chapters~\ref{ch:candidates} through~\ref{ch:arch} are beside the point.

\begin{table}[t]
\centering
\small
\begin{tabular}{lrrr}
\toprule
\textbf{System} & \textbf{Precision} & \textbf{Recall} & \textbf{MCC} \\
\midrule
GPT-4o, zero-shot        & \textbf{8.45\%} & 50.0\% & $+0.109$ \\
GPT-4o, few-shot         & 4.96\%          & 58.3\% & $+0.012$ \\
ACCORD, $\tau = 0.5$     & 4.24\%          & 77.8\% & $-0.007^{\dagger}$ \\
Llama 3 8B (evaluation baseline) & 4.55\%        & 91.7\% & $-0.044$ \\
\bottomrule
\end{tabular}
\caption{System comparison on the $253$ non-ambiguous pharmacist-annotated
pairs ($12$ positive, $241$ negative). Matthews
correlation coefficient \citep{matthews1975mcc} is reported because it is the
summary statistic least distorted by extreme class imbalance; a value of $0$ is
chance. ACCORD's precision and recall are the decontaminated three-seed means
($4.24\% \pm 0.23$, $77.8\% \pm 3.9$); its MCC, marked $\dagger$, is computed
from the \emph{contaminated} confusion matrix on the $253$-pair core and has not
been recomputed under decontamination. \textbf{The Llama row is a separately
prompted Llama 3\,8B evaluation run on these pairs, using the original
labeller's prompt, and is not the weak-supervision training labels
themselves}; measured directly against its own labels, the weak labeller
scores MCC $+0.024$, precision $4.93\%$, recall $91.7\%$ (Chapter~\ref{ch:weak}).
All four MCCs carry $95\%$ bootstrap
intervals that contain zero (Chapter~\ref{ch:metrics},
Table~\ref{tab:met-mccci}), so the ordering they induce over these four systems
is not identified and no claim in this thesis rests on it. What the table
supports is weaker and sufficient: no system here is distinguishable from chance
on this evaluation set, and the two highest-recall systems achieve their recall
by predicting the positive class almost everywhere.}
\label{tab:systems}
\end{table}

\subsection{The four findings in Table~\ref{tab:systems}}

\paragraph{The task is not impossible.} GPT-4o zero-shot
\citep{openai2023gpt4} achieves $8.45\%$ precision at $50\%$ recall with an MCC
of $+0.109$. This is not a usable clinical system, and it should not be presented
as one: eleven false alarms per true contradiction, and half the true
contradictions missed. Nor is it established to be above chance: its $95\%$
bootstrap interval is $[-0.0319, 0.2484]$ and contains zero
(Chapter~\ref{ch:metrics}, Table~\ref{tab:met-mccci}). The claim that survives is
the narrower one that its precision of $8.45\%$ is roughly double the base rate
on the same evidence and the same $253$ pairs, at a predicted-positive rate far
below one. [Correction note: the earlier text read ``it is above chance,'' which
the interval does not support; whether the evidence contains signal is therefore
undetermined at this sample size rather than settled in either direction.]

\paragraph{Few-shot prompting made GPT-4o worse.} Precision falls from $8.45\%$
to $4.96\%$ and MCC from $+0.109$ to $+0.012$, while recall rises from $50.0\%$
to $58.3\%$. The exemplars pulled the model towards the positive class. This is
consistent with the class-prior sensitivity documented for in-context learning
\citep{zhao2021calibrate}, and it is the same failure direction as the weak
labeller's, arrived at independently. [TODO: report the exemplar class balance
used in the few-shot prompt. If it over-represents positives relative to the
$4\%$ base rate, the mechanism is identified rather than hypothesised.]

\paragraph{ACCORD did not improve on its teacher.} ACCORD ($-0.007$) sits below
the weak labeller's own MCC on these pairs, $+0.024$ (TP~$11$, FN~$1$,
FP~$212$, TN~$29$) --- both are point estimates indistinguishable from chance,
not evidence that either scores below it. [Correction note: an earlier draft
put the teacher's MCC at $-0.044$ here, which is the score of a separately
prompted Llama 3\,8B evaluation run (the Llama row of
Table~\ref{tab:systems}), not the weak labeller's own calls.] Under
decontamination ACCORD's recall falls from the teacher's $91.7\%$ to
$77.8\% \pm 3.9$ while its precision, $4.24\%$, drops below both the teacher's
$4.93\%$ and the $4.74\%$ base rate. Both remain near-constant positive
predictors. This is Chapter~\ref{ch:weak}'s finding restated on the gold set: the
student reproduced the teacher's marginal, and reproduced it slightly worse.

\paragraph{The purpose-built system is not distinguishable from any of the
other three.} A model carrying roughly $7.5$M task-specific parameters over a
domain-pretrained biomedical encoder, trained on $1{,}679$ in-domain examples,
does not beat a general-purpose model given a single prompt and no training data
at all, and does not beat the $8$B labeller it was distilled from. What cannot be
said is that it is ranked third or fourth of four: all four bootstrap intervals
overlap and all four contain zero. [Correction note: the earlier heading read
``the second-worst of four'' and the earlier text ``performs worse on the target
task,'' both of which rely on an ordering that
Chapter~\ref{ch:metrics}, Table~\ref{tab:met-mccci}, shows is not identified.]

\section{Error Taxonomy}
\label{sec:errors}

\subsection{What the twelve true contradictions are}

\begin{table}[t]
\centering
\small
\begin{tabular}{lrr}
\toprule
\textbf{Contradiction type} & \textbf{Count} & \textbf{Share} \\
\midrule
Omission (drug present in one note, absent from the other) & 7 & 58.3\% \\
Dosage discrepancy                                          & 3 & 25.0\% \\
Drug substitution                                           & 1 & 8.3\% \\
Class conflict                                              & 1 & 8.3\% \\
\midrule
\textbf{Total}                                              & \textbf{12} & \\
\bottomrule
\end{tabular}
\caption{Types of the twelve pharmacist-confirmed contradictions. Omission
dominates. This matches the medication reconciliation literature, where omission
of a regularly used medication is the single most common category of unintended
discrepancy \citep{cornish2005discrepancies}, and it means the hardest evidential
case, reasoning about something that is \emph{not} in a document, is also the
most clinically important one.}
\label{tab:errtypes}
\end{table}

\begin{table}[t]
\centering
\small
\begin{tabular}{lrrrrr}
\toprule
\textbf{Danger rating} & \textbf{1} & \textbf{2} & \textbf{3} & \textbf{4} & \textbf{5} \\
\midrule
Count & 2 & 4 & 2 & 1 & 3 \\
\bottomrule
\end{tabular}
\caption{Pharmacist danger ratings of the twelve true contradictions. Four of
twelve are rated $4$ or $5$. These are the cases a triage system exists to
surface, and Section~\ref{sec:m3} established that no model score orders them.}
\label{tab:danger}
\end{table}

\subsection{Who catches what}

\begin{table}[h]
\centering
\small
\begin{tabular}{lr}
\toprule
\textbf{System} & \textbf{True contradictions recovered} \\
\midrule
Llama 3 8B          & 11 / 12 \\
ACCORD              & 9.3 / 12 (seeds: 9, 10, 9) \\
GPT-4o, zero-shot   & 6 / 12 \\
\bottomrule
\end{tabular}
\caption{Recall on the twelve true contradictions. The ACCORD row is the
decontaminated configuration's three seeds; the contaminated run reported
$11/12$, which no decontaminated seed reaches. The two near-constant positive
predictors recover most of the twelve; the selective predictor recovers six. This
table should not be read as evidence that the first two are better detectors,
and Section~\ref{sec:errors-reading} explains why.}
\label{tab:recall}
\end{table}

\subsection{How to read Table~\ref{tab:recall}}
\label{sec:errors-reading}

Recall in isolation is the most misleading metric available in this setting. A
system that predicts \emph{contradiction} for everything achieves recall
$1.0$. ACCORD and Llama recover most of the twelve because they predict positive
for the overwhelming majority of pairs: ACCORD for $278$ of $299$ in the
contaminated run$^{\dagger}$, and Llama at a comparable rate, consistent with the
$84.9\%$ positive rate Chapter~\ref{ch:weak} measured over the full candidate
pool. GPT-4o achieves $6/12$ while predicting positive far less often, and is the
only system whose MCC point estimate is meaningfully positive --- though its
interval, like the other three, contains zero.

The one substantive observation available from Table~\ref{tab:recall} concerns
where ACCORD's misses fall. Chapter~\ref{ch:candidates}, Section~\ref{sec:gold},
established arithmetically that ACCORD recovers one of the two true Rule 4
contradictions, and Rule 4 is the only rule on which the model exercises
meaningful selectivity ($74.5\%$ positive rather than $\approx 97\%$). The one
place the model declines to predict positive is the one place it declines
incorrectly. There is no configuration of this model in which its selectivity
helps. [Correction note: the per-rule recovery arithmetic in
Chapter~\ref{ch:candidates} was derived from the contaminated run, which missed
one of twelve; the decontaminated seeds miss two or three. The per-rule
attribution of those additional misses is not recorded in the decontaminated
artifacts and is not recomputed here.]

\subsection{The clinically important observation}

Omission is $7$ of $12$ true contradictions and it is the category the
architecture is least equipped for. Detecting an omission requires reasoning
about absence, and the model's input for a Rule 4 candidate is a nursing
sentence paired with an empty string. There is no discharge evidence to encode,
so an asymmetric cross-note encoder has nothing to attend to on one side, and
the pair representation degenerates to a single-sentence representation. A
system designed around cross-note attention is structurally poorly matched to
the majority error type in its own target population. This is a design finding
independent of every implementation defect in Chapter~\ref{ch:arch}, and it
would apply equally to a correct implementation of the same specification.

\section{What the Exemplar Defect Does and Does Not Contaminate}
\label{sec:eval-scoping}

The most serious objection to this chapter is that Chapter~\ref{ch:candidates}
established that $41.4\%$ of action-conflict candidates were annotated against
evidence that did not carry their own conflict. If the labels are contaminated,
are the measurements meaningless? The answer requires separating three distinct
claims.

\subsection{What is genuinely bounded}

The claim ``the candidate pool has a precision of $4.0\%$'' is an
\emph{upper bound on the validity of the annotation exercise}, not an estimate
of pipeline quality. The pharmacist rendered correct judgements about the
evidence presented; for a large minority of action-conflict candidates that
evidence was not the evidence the rule had found. We do not know how the same
annotator would have ruled on corrected exemplars, and both a higher and a
lower true precision are consistent with the data collected. Every use of the
$4.0\%$ figure in this thesis carries this qualification.

\subsection{What is not contaminated}

The four measurements in this chapter are statements about whether the model's
outputs carry information about the pharmacist's labels. Both sides of each
comparison refer to \emph{the same object}: the model scored exactly the
exemplar sentence pairs the pharmacist annotated. The question ``does the model's
score order the pharmacist's label'' is therefore well-posed regardless of
whether that label is a good proxy for clinical truth. If the model had learned
any function of the presented sentence pair, and the pharmacist's judgement is
any function of that same pair, the two would correlate. They do not.

The same argument covers the threshold sweep, the severity correlations, the
weak-label marginal of Chapter~\ref{ch:weak}, and the transfer collapse of
Chapter~\ref{ch:arch}. None of these compares a model output against a claim
about clinical reality; all compare model outputs against each other or against
annotations of the identical input.

\subsection{The residual risk, stated honestly}

There is a real residual risk and it should not be minimised. Label noise
attenuates a measured AUROC towards $0.5$. In principle, a model with a true
AUROC of, say, $0.65$ could be measured at $0.5145$ under sufficient label noise.
Four considerations bound how much weight this deserves.

\begin{enumerate}
\item \textbf{Most of the evaluation set is unaffected.} The defect applies only
      to action-conflict candidates, Rules 1 and 2, which are $165$ of the $300$
      sampled pairs. The direct measurement on these pairs (\texttt{exemplar\_correctness\_recomputed.json}) gives $65$ of $137$ resolved rule-1 and rule-2 pairs wrong, rather than the $68$ an applied pooled rate would estimate;
      affected pairs, under a quarter of the evaluation set. Rules 3 and 4,
      $134$ of the $299$ annotated pairs, are free of the defect by
      construction, because Rule 3
      selects the maximum-difference dose pair and Rule 4 has no discharge
      exemplar to mis-select.
\item \textbf{The decisive experiment has since been run.}
      Chapter~\ref{ch:discussion}, Section~\ref{sec:disc-strata}, reports the
      evaluation restricted to the Rule 3 and Rule 4 strata, which are free of
      the exemplar defect: AUROC $0.5078$ at Mann-Whitney $p = 0.953$, with a
      bootstrap $95\%$ interval of $[0.2217, 0.7615]$ over $B = 10{,}000$
      resamples. Two properties of that interval must be reported with it. Of
      the $10{,}000$ resamples, $20$ ($0.20\%$) were dropped because they
      contained no positives and AUROC is undefined on them; and the stratified
      interval, which preserves the $6/128$ class split in every resample, is
      $[0.2643, 0.7344]$. An interval that silently discards undefined
      resamples is the same class of omission this chapter is about, so the
      count is stated rather than absorbed. [TODO: report average precision on
      the restricted strata against the restricted base rate.]
\item \textbf{A near-constant predictor has no signal to attenuate.} The model
      predicts positive for $278$ of $299$ pairs and all its scores lie above
      $0.15$ in a narrow band. Label noise degrades an ordering; it does not
      explain the near-absence of an ordering, and Chapter~\ref{ch:weak}
      supplies an independent, sufficient explanation for that absence in the
      $21\times$ prevalence mismatch of the supervision.
\item \textbf{The top-$k$ result is robust to moderate label noise.} A model
      with genuine signal concentrates true positives at the top of its ranking
      even under substantial label corruption. One true contradiction in the top
      fifteen, against the $0.6$ expected under random ordering, is not what an
      attenuated signal looks like.
\end{enumerate}

Point 2 was a genuine hole in the earlier version of this work; it is now closed
by measurement, at six positives and with the interval width that implies.

\section{Synthesis}
\label{sec:synthesis}

\begin{table}[t]
\centering
\small
\begin{tabular}{p{0.3\linewidth}p{0.2\linewidth}p{0.42\linewidth}}
\toprule
\textbf{Stage} & \textbf{Status} & \textbf{Decisive measurement} \\
\midrule
Cohort construction & Works & $9{,}783$ admissions, $19{,}566$ notes, documented
retreat from three roles to two \\
\addlinespace
Medication extraction & Unvalidated & $8.9{:}1$ discharge/nursing asymmetry;
never annotated \\
\addlinespace
Candidate generation & Enriches, but evidence selection was wrong & Old exemplar
carried the conflict $58.6\%$ of the time; $739/1{,}307$ changed \\
\addlinespace
Weak supervision & Fails & $84.9\%$ positive against a $4.0\%$ base rate;
$21\times$ \\
\addlinespace
Architecture & Three claims falsified & Authority parameter moved $1.8\%$/$0.6\%$
from its initialisation; rule $F_1$ $0.96 \to 0.065$ under transfer; harm head
absent from the loss \\
\addlinespace
Detection & Fails & AUROC $0.5145$, CI $[0.321, 0.703]$, $p = 0.867$; one true
positive in the top fifteen against $0.6$ expected \\
\addlinespace
Severity ranking & Not tested & $\rho = -0.0001$ from an untrained head \\
\bottomrule
\end{tabular}
\caption{Where the pipeline stands, stage by stage. Two stages work, one is
unvalidated, one is defective, and three fail. The failures are independent:
fixing the exemplar selector would not correct the supervision marginal, and
correcting the marginal would not connect the harm head to the loss.}
\label{tab:synthesis}
\end{table}

\subsection{The claim this thesis makes}

The standard pipeline for this problem, rule-based candidate mining, language
model weak supervision, and a purpose-built neural detector, fails at every
stage, and the failures are measurable and diagnosable rather than incidental.
Table~\ref{tab:synthesis} is that claim in full.

The claim is about this pipeline, not about the construction in general, and
the two are not interchangeable. The base rate the rules produced, the
labeller's over-calling at that base rate, and the twelve-positive evaluation
are properties of the construction: any group applying high-recall rules,
LLM weak supervision, and a fixed-size expert sample to an event this rare
should expect something in the same range. The exemplar selector's
positional indexing (Section~\ref{sec:exemplar}) and the harm head's absent
loss term (Section~\ref{sec:c3}) are properties of this build. A
differently written selector or a correctly wired loss would not repeat
either one, and neither defect is evidence that rule-based mining or
guideline-grounded severity scoring cannot be done correctly.

Three properties of the failure make it worth reporting rather than repairing
quietly.

\paragraph{It is not one failure.} The exemplar defect, the prevalence mismatch,
the initialisation artifact, the label leakage, and the untrained head are
causally independent. Each has its own mechanism and its own detector. A project
that fixed any one of them would still have produced a system with AUROC near
$0.5$.

\paragraph{Every failure was detectable before the results table.} Not one of
them required a held-out set, an expensive annotation, or a novel diagnostic.
The full detection kit is: a base-rate estimate before weak labelling, a printed
label marginal, a distinct-value count on the confidence field, a module
manifest with loss-graph membership, a parameter-displacement log, an
out-of-domain transfer run, an output-range check, a majority-class baseline,
and a group-disjointness assertion. Chapter~\ref{ch:arch},
Section~\ref{sec:diagnostics} states them as a protocol.

\paragraph{The negative result is measured four ways.} Ranking, thresholds,
severity, and comparison against three external systems. They fail for different
reasons if they fail, and they agree.

\subsection{What remains true after all of this}

The clinical problem is real. Omission of a regularly used medication is the
most common category of unintended discrepancy at care transitions
\citep{cornish2005discrepancies}, and adverse events following discharge are
common, frequently drug-related, and frequently preventable
\citep{forster2003adverse}. A board-certified pharmacist found $12$ genuine
contradictions in $299$ pairs, four of them rated $4$ or $5$ for danger. GPT-4o
zero-shot detects half of them above chance. Nothing in this thesis suggests the
problem is not worth solving. It establishes that this particular route to
solving it does not work, and it identifies, at each stage, the measurement that
would have said so earlier.

\section{Limitations}
\label{sec:eval-limitations}

\paragraph{Twelve positives.} Every confidence interval in this chapter is wide
because the evaluation set contains twelve positive examples. The bootstrap
interval on AUROC, $[0.321, 0.703]$, is consistent with a genuinely inverted
model and with a mildly informative one. What it is \emph{not} consistent with
is a useful one, and that is the claim being made. A larger annotated set is the
single most valuable extension of this work.

\paragraph{Neither configuration is subject-disjoint, and the quantities not
recomputed.} Section~\ref{sec:m1} states the leakage position in full: the
contaminated run leaked at row level, the decontaminated replacement leaks at
subject level ($158$ of $266$ evaluated subjects appear in its training pool),
and only the C19 control at $0.497 \pm 0.035$ is certified subject-disjoint. A
second, narrower gap follows from the same correction. Several quantities in
this chapter derive from the contaminated run and have not been recomputed under
decontamination; they are marked $\dagger$ where they appear. They are: the
predicted-positive count of $278$ and the accuracy of $0.115$ in
Table~\ref{tab:point} and the denominator reconstruction that rests on them; the
ambiguous-as-positive and harm-score rows of Table~\ref{tab:auroc}; the
threshold sweep and the precision-at-fixed-recall table of Section~\ref{sec:m2};
the Spearman correlations of Section~\ref{sec:m3}; and ACCORD's MCC in
Table~\ref{tab:systems}. [TODO: recompute each of these from
\texttt{gold\_eval\_decontam\_predictions.csv} and remove the daggers.]

\paragraph{One annotator, no inter-rater reliability.}
Chapter~\ref{ch:candidates}, Section~\ref{sec:gold}. The previously reported
$\kappa = 1.000$ is retracted as circular. Inter-rater reliability is not
established for any label used in this chapter.

\paragraph{The exemplar defect.} Section~\ref{sec:eval-scoping}, including the
[TODO] that would close the objection.

\paragraph{Recall is measured only within the candidate pool.} Nothing outside
the $1{,}978$ candidates is detectable by any system evaluated here, including
GPT-4o, which was given the same candidates. The rules' recall against all true
contradictions in the $9{,}783$-admission cohort is unmeasured.

\paragraph{No efficiency analysis is reported.} Training cost, inference
latency, and the cost of the GPT-4o comparison were not systematically recorded.
[TODO: report wall-clock training time, inference latency per pair, and total
API cost for the GPT-4o rows. For a system whose value proposition is triage at
scale, cost per screened admission is a first-order metric and its absence is a
gap.]

\paragraph{One backbone, one adapter configuration.} Every ACCORD result uses
BioMedBERT with LoRA at $r = 16$. Section~\ref{sec:mednli} shows this
configuration costs eleven accuracy points against full fine-tuning on MedNLI.
Whether a full fine-tune of the same architecture would behave differently is
untested. It would not repair the supervision, but the claim that it changes
nothing is not established. [TODO: retrain ACCORD with a full backbone
fine-tune on the identical held-out split and report AUROC.]

\paragraph{Single institution.} MIMIC-III only. The transfer results in
Chapter~\ref{ch:arch} move across index diagnosis, not across institution.

\section{Ethics Statement}
\label{sec:eval-ethics}

\paragraph{Data governance.} MIMIC-III is credentialed data used under a data
use agreement. No note text appears in this thesis or in any derived artifact,
all note-derived intermediates are confined to a gitignored local directory, and
no re-identification was attempted. The one open governance question concerns
the GPT-4o comparison in Section~\ref{sec:m4}. [TODO: state the governance basis
under which candidate excerpts were transmitted to a hosted API, including
whether zero-retention terms were in force. If this cannot be established, the
GPT-4o rows should carry that caveat wherever they are reported.]

\paragraph{Deployment risk of the system as measured.} A system with $77.8\%$
recall and $4.24\%$ precision deployed as a clinical alert generates more than
twenty false alarms per true contradiction. The established consequence of alert
burdens at that ratio is desensitisation, and desensitisation degrades response
to the alerts that matter. \textbf{This system should not be deployed in any
form.} The stronger statement is that its failure mode is worse than doing
nothing, because it would consume clinician attention while providing, by
Section~\ref{sec:m1}, no information.

\paragraph{Automation bias.} A component of this work presented a severity score
to a downstream consumer. Chapter~\ref{ch:arch} established that the score comes
from an untrained head. A clinical decision support system displaying a
plausible-looking numeric danger score derived from random initialisation is a
specific and serious hazard, because numeric scores in clinical interfaces are
acted upon. That this reached a results table is the most serious individual
failure documented in this thesis, and the diagnostic that catches it, printing
the empirical range of every emitted score, costs one line.

\paragraph{Publication ethics.} Four claims that appeared in an earlier
submission of this work are retracted here: the discovered authority hierarchy,
the rule-type $F_1$ of $0.9603$, the internal-consistency $\kappa$ of $1.000$,
and the intermediate-documentation prefilter finding. Two further reporting
defects are corrected. First, and most seriously, the earlier version of this
thesis stated four separate times that the configuration trained on the
evaluation set was withdrawn and reported nowhere; in fact it was the source of
every number in Chapters~\ref{ch:eval} and~\ref{ch:metrics} and of the
defect-free analysis in Chapter~\ref{ch:discussion}. Those figures are replaced
throughout by the decontaminated retrain, which is worse on recall, precision,
$F_1$ and $\kappa$, so the correction strengthens the negative conclusion; the
false statements of non-reporting are themselves retracted here. Second, an
ablation table reporting three seeds is corrected to
two seeds plus a hardcoded reference row. Retracting one's own claims in the
document that supersedes them is the minimum standard, not a virtue, and it is
recorded here so that any reader encountering the earlier artifacts knows their
status.

\paragraph{Environmental cost.} Not measured. [TODO: report GPU-hours for
training and for the transfer runs, and estimate the associated energy use.]

\paragraph{Dual use and population bias.} The cohort is critical-care patients
at one US academic medical centre in a historical period. The drug dictionary,
the class dictionary, and the danger judgements encode that context. Any system
built on this pipeline would inherit those biases, and its errors would not be
distributed uniformly across patient populations. Nothing in this work measured
subgroup performance, and with twelve positives nothing could.

\section{Summary}
\label{sec:eval-summary}

BioMedBERT fine-tuned on MedNLI reaches $83.90\%$ accuracy and $83.91\%$ macro
$F_1$, inside the published range. That result is a property of the backbone and
carries no ACCORD components; reporting it under the system's name would be the
most misleading available framing and is avoided.

On the $299$ pharmacist-annotated pairs, ACCORD predicts \emph{contradiction} for
the large majority, achieving precision $0.0424 \pm 0.0023$ against a base rate
of $0.040$ and recall $0.7778 \pm 0.0393$. These are the decontaminated
configuration's three-seed values; the corresponding figures in earlier drafts
($0.047$ precision, $0.917$ recall) came from the contaminated run. Four
independent measurements establish that this is not a threshold problem:

\begin{enumerate}
\item \textbf{Ranking.} AUROC $0.5145$ for the three-seed ensemble
      ($0.5211 \pm 0.0201$ per seed), average precision $0.0579 \pm 0.0034$
      against a $0.040$ base rate, Mann-Whitney $p = 0.867$, bootstrap $95\%$ CI
      $[0.321, 0.703]$. The fifteen highest-scored pairs contain one true
      contradiction against $0.6$ expected under random ordering. Neither this
      configuration nor the contaminated one it replaces is certified
      subject-disjoint; the only subject-disjoint measurement in the project is
      the plain control at $0.497 \pm 0.035$.
\item \textbf{Operating points.} At thresholds $0.05$, $0.10$, and $0.15$ the
      model predicts positive for every pair at precision $0.0474$, exactly the
      base rate. Precision at recall $0.5$ is $5.4\%$.
\item \textbf{Severity.} Spearman $\rho = -0.0001$ ($p = 0.999$) between harm
      score and danger rating, from a head that never received a gradient. The
      severity research question was not tested.
\item \textbf{Comparison.} GPT-4o zero-shot is the only system whose MCC point
      estimate is meaningfully positive ($+0.109$); ACCORD sits at $-0.007$ and
      the weak labeller it was distilled from at $+0.024$, both indistinguishable
      from chance. [Correction note: an earlier draft put the weak labeller's
      score at $-0.044$ here, which is the score of a separately prompted Llama
      3\,8B evaluation run, not the weak labeller's own calls; that run, not the
      weak labeller, is the fourth row of Table~\ref{tab:met-mccci}.] All four
      bootstrap intervals in that table contain zero (Chapter~\ref{ch:metrics}),
      so no system in the comparison is distinguishable from chance and the
      ordering over them is not identified.
\end{enumerate}

Omission accounts for $7$ of the $12$ true contradictions, and it is the error
type an asymmetric cross-note encoder is structurally least able to represent,
because one side of the input is empty by construction. Four of the twelve are
rated $4$ or $5$ for danger, and no model score orders them.

Whether the task is possible is not settled here: a general-purpose model,
zero-shot, on the same evidence, reaches roughly double the base rate in
precision while predicting positive far less often, but its MCC interval contains
zero at this sample size. What is settled is that this pipeline fails, and that
it fails at every stage, for independent and individually diagnosable reasons.
